\documentclass[11pt,a4paper]{article}
\usepackage{amsmath}
\usepackage{amssymb}
\usepackage{geometry}
\usepackage{hyperref}

\geometry{margin=2.5cm}

\title{Numerical Methodology for Computing Poloidal Spectrograms}
\author{}
\date{\today}

\begin{document}

\maketitle

\section{Overview}

This document describes the numerical procedure for computing frequency-wavenumber spectrograms of the electrostatic potential $\phi$ in the poloidal direction. The spectrogram reveals the dispersion relation $\omega(k_\theta)$ of fluctuations at a fixed radial location.

\section{Data Extraction}

\subsection{Poloidal Cut at Fixed Radius}

At each time frame $t_i$ (with $i = 1, \ldots, N_t$), the electrostatic potential is extracted along a poloidal cut at radial position $x_0$:
\begin{equation}
\phi(\theta, t_i) \quad \text{at} \quad x = x_0
\end{equation}

The poloidal coordinate $\theta$ is discretized with $N_\theta$ points, and the arclength along the poloidal direction is computed as:
\begin{equation}
\ell_j = \sum_{k=1}^{j} \sqrt{(\Delta R_k)^2 + (\Delta Z_k)^2}
\end{equation}
where $(R, Z)$ are the cylindrical coordinates in the poloidal plane and $\Delta R_k$, $\Delta Z_k$ are the differences between adjacent $k$-th grid points.

\subsection{Poloidal ExB Velocity}

\subsubsection{Magnetic Field Decomposition}

In flux coordinates, the magnetic field is decomposed into toroidal and poloidal components:
\begin{equation}
\mathbf{B} = B_\phi \, \mathbf{e}_\phi + B_\theta \, \mathbf{e}_\theta
\end{equation}

We introduce the pitch angle factor $\alpha$ to characterize the relative strength of these components:
\begin{equation}
\alpha = \frac{1}{1 + \frac{r}{Rq}}
\end{equation}
where $r$ is the minor radius, $R$ is the major radius, and $q$ is the safety factor. This allows us to approximate:
\begin{equation}
\mathbf{B} \approx |\mathbf{B}| \left( \alpha \, \mathbf{e}_\phi + (1 - \alpha) \, \mathbf{e}_\theta \right)
\end{equation}

\subsubsection{ExB Velocity Derivation}

The $\mathbf{E} \times \mathbf{B}$ drift velocity is given by:
\begin{equation}
\mathbf{v}_E = \frac{\mathbf{E} \times \mathbf{B}}{B^2} = -\frac{\nabla \phi \times \mathbf{B}}{B^2} = -\frac{1}{|\mathbf{B}|} \nabla \phi \times \frac{\mathbf{B}}{|\mathbf{B}|}
\end{equation}

Substituting the magnetic field decomposition:
\begin{equation}
\mathbf{v}_E = -\frac{1}{|\mathbf{B}|} \nabla \phi \times \left( \alpha \, \mathbf{e}_\phi + (1 - \alpha) \, \mathbf{e}_\theta \right)
\end{equation}

Expanding the gradient in local flux coordinates $(r, \theta, \phi)$:
\begin{equation}
\nabla \phi = \frac{\partial \phi}{\partial r} \mathbf{e}_r + \frac{1}{r} \frac{\partial \phi}{\partial \theta} \mathbf{e}_\theta + \frac{1}{R} \frac{\partial \phi}{\partial \phi} \mathbf{e}_\phi
\end{equation}

For axisymmetric fluctuations ($\partial/\partial \phi = 0$), the cross products yield:
\begin{align}
\mathbf{v}_E &= -\frac{1}{|\mathbf{B}|} \left[ \alpha \left( \frac{\partial \phi}{\partial r} \mathbf{e}_r \times \mathbf{e}_\phi + \frac{1}{r} \frac{\partial \phi}{\partial \theta} \mathbf{e}_\theta \times \mathbf{e}_\phi \right) \right. \nonumber \\
&\quad \left. + (1 - \alpha) \left( \frac{\partial \phi}{\partial r} \mathbf{e}_r \times \mathbf{e}_\theta + \frac{1}{r} \frac{\partial \phi}{\partial \theta} \mathbf{e}_\theta \times \mathbf{e}_\theta \right) \right]
\end{align}

Using the relations $\mathbf{e}_r \times \mathbf{e}_\phi = -\mathbf{e}_\theta$, $\mathbf{e}_\theta \times \mathbf{e}_\phi = \mathbf{e}_r$, $\mathbf{e}_r \times \mathbf{e}_\theta = \mathbf{e}_\phi$, and $\mathbf{e}_\theta \times \mathbf{e}_\theta = 0$:
\begin{equation}
\mathbf{v}_E = -\frac{1}{|\mathbf{B}|} \left[ \alpha \left( \frac{\partial \phi}{\partial r} (-\mathbf{e}_\theta) + \frac{1}{r} \frac{\partial \phi}{\partial \theta} \mathbf{e}_r \right) + (1 - \alpha) \frac{\partial \phi}{\partial r} \mathbf{e}_\phi \right]
\end{equation}

Collecting components:
\begin{align}
v_{E,r} &= -\frac{\alpha}{|\mathbf{B}|} \frac{1}{r} \frac{\partial \phi}{\partial \theta} \\
v_{E,\theta} &= \frac{\alpha}{|\mathbf{B}|} \frac{\partial \phi}{\partial r} \\
v_{E,\phi} &= -\frac{(1 - \alpha)}{|\mathbf{B}|} \frac{\partial \phi}{\partial r}
\end{align}

\subsubsection{Numerical Implementation}

The poloidal component is computed at each point $(r, \theta, t)$ using centered finite differences for the radial derivative:
\begin{equation}
v_{E,\theta}(r, \theta, t) = \frac{\alpha}{B(r,\theta)} \frac{\phi(r + \Delta r, \theta, t) - \phi(r - \Delta r, \theta, t)}{2 \Delta r}
\end{equation}

The time- and poloidal-averaged poloidal rotation velocity is then:
\begin{equation}
\langle v_{E,\theta} \rangle_{t,\theta} = \frac{1}{N_t N_\theta} \sum_{i=1}^{N_t} \sum_{j=1}^{N_\theta} v_{E,\theta}(\theta_j, t_i)
\end{equation}

\section{Spectral Analysis}

\subsection{Windowing in Time}

To enforce periodicity in time and reduce spectral leakage, a Kaiser window with parameter $\beta = 8$ is applied:
\begin{equation}
\tilde{\phi}(\theta, t_i) = \phi(\theta, t_i) \cdot w_{\text{Kaiser}}(t_i)
\end{equation}

\subsection{Poloidal Fourier Transform}

For each time frame, the real FFT is computed in the poloidal direction:
\begin{equation}
\hat{\phi}(k_\theta, t_i) = \text{FFT}_\theta[\tilde{\phi}(\theta, t_i)]
\end{equation}

The poloidal wavenumber range is:
\begin{equation}
k_\theta \in \left[\frac{2\pi}{L}, \frac{N_\theta \pi}{L}\right]
\end{equation}
where $L = \ell_{N_\theta}$ is the total poloidal circumference.

\subsection{Doppler Shift Correction}

To account for the poloidal rotation due to $\mathbf{E} \times \mathbf{B}$ drift, a phase correction is applied in the laboratory frame:
\begin{equation}
\hat{\phi}_{\text{shift}}(k_\theta, t_i) = \hat{\phi}(k_\theta, t_i) \cdot \exp\left(-i k_\theta \langle v_{E,\theta} \rangle_{t,\theta} t_i\right)
\end{equation}

This transformation moves to a frame rotating with the mean poloidal ExB velocity, removing the Doppler shift:
\begin{equation}
\omega_{\text{lab}} = \omega_{\text{plasma}} + k_\theta \langle v_{E,\theta} \rangle_{t,\theta}
\end{equation}

\subsection{Temporal Fourier Transform}

The frequency spectrum is obtained by Fourier transforming in time:
\begin{equation}
\hat{\phi}(k_\theta, \omega) = \text{FFT}_t[\hat{\phi}_{\text{shift}}(k_\theta, t)]
\end{equation}

The result is shifted to center zero frequency:
\begin{equation}
\hat{\phi}(k_\theta, \omega) = \text{fftshift}\left[\text{FFT}_t[\hat{\phi}_{\text{shift}}(k_\theta, t)]\right]
\end{equation}

The frequency range is:
\begin{equation}
\omega \in \left[-\frac{N_t \pi}{T}, \frac{N_t \pi}{T}\right]
\end{equation}
where $T = t_{N_t} - t_1$ is the total simulation time analyzed.

\section{Normalization and Visualization}

\subsection{Normalized Units}

The spectrogram is presented in gyrokinetic normalized units:
\begin{align}
k_\theta^* &= k_\theta \rho_s \\
\omega^* &= \frac{\omega R}{c_s}
\end{align}
where:
\begin{itemize}
    \item $\rho_s = c_s / \Omega_i$ is the ion sound Larmor radius
    \item $c_s = \sqrt{T_e / m_i}$ is the ion sound speed
    \item $\Omega_i = eB / m_i$ is the ion cyclotron frequency
    \item $R$ is the major radius
\end{itemize}

\subsection{Power Spectral Density}

The spectrogram displays the power spectral density:
\begin{equation}
S(k_\theta, \omega) = |\hat{\phi}(k_\theta, \omega)|^2
\end{equation}

This is visualized using a logarithmic color scale to capture the wide dynamic range of fluctuation amplitudes.

\section{Physical Interpretation}

The spectrogram $S(k_\theta, \omega)$ reveals:
\begin{itemize}
    \item The dispersion relation $\omega(k_\theta)$ of dominant modes
    \item The spectral width indicating coherence and damping
    \item The propagation direction from the sign of $\omega/k_\theta$
    \item Mode identification (ITG, TEM, ETG) from characteristic slopes
\end{itemize}

After Doppler shift correction, positive frequencies indicate propagation in the electron diamagnetic direction, while negative frequencies indicate ion diamagnetic direction, both in the plasma frame.

\end{document}
